\documentclass[10pt,a4paper]{article}
\usepackage{fontspec,amsmath,amssymb,graphicx}
\usepackage[margin=18mm]{geometry}
\setmainfont{Times New Roman}
\setlength{\parindent}{0pt}
\setlength{\parskip}{5pt}
\allowdisplaybreaks
\sloppy
\emergencystretch=2em
\begin{document}
\section*{BT\_5\_6\_3\_07 — Lời giải bổ sung}
Đặt $\mathcal H$ là phép kéo dài điều hòa bốn cạnh trong hình chữ nhật theo công thức của bài 5.1.3.1, và $\mathcal H_N$ là phép kéo dài với hai cạnh đứng Neumann 0 (chuỗi cosine theo $x$, sinh theo $y$). Một nghiệm riêng $P$ cho nguồn giúp giải ngay $u=P+$phép kéo dài của dữ liệu biên trừ $P$.

(a) $P=-\sin2\pi x/(4\pi^2)+\sin3\pi x/(3\pi^2)$; $u=P+\mathcal H[-P,-P,0,0]$ trên hình vuông đơn vị.
(b) $P=-\sin\pi y/\pi^2$; $u=P+\mathcal H[-1,1,\sin\pi y/\pi^2,\sin\pi y/\pi^2]$.
(c) $P=-\sin2x/4-\sin3y/9$; $u=P+\mathcal H[-1+\sin2x/4,1+\sin2x/4,-1+\sin3y/9,1+\sin3y/9]$ trên $(0,\pi)^2$.
(d) $P=x^3y/6$; $u=P+\mathcal H[0,1-\pi x^3/6,0,-\pi^3y/6]$.
(e) Trên $(0,\pi)^2$, lấy cơ sở $X_n=\cos nx$, biên dưới 0, trên 1. Nguồn $F=xy$ có $F_0(y)=\pi y/2$, $F_n(y)=2[(-1)^n-1]y/(\pi n^2)$, $n\ge1$. Mode 0 $U_0=y/\pi+\pi(y^3-\pi^2y)/12$. Các mode $n\ge1$:
\[U_n(y)=\int_0^\pi\frac{-\sinh(n\min(y,s))\sinh(n[\pi-\max(y,s)])}{n\sinh n\pi}F_n(s)ds.\]
$u=U_0+\sum U_n\cos nx$. Biên Neumann theo $x$ thỏa tự động.\end{document}